\textbf{Phần 1. Trắc nghiệm lựa chọn}
\setcounterpageref{ex}{0}

\begin{ex} %[5H9N49]
Hình khai triển của một hình lập phương gồm bao nhiêu hình vuông bằng nhau?
\choice
{$4$ hình vuông}
{$8$ hình vuông}
{\True $6$ hình vuông}
{$12$ hình vuông}
\loigiai{
Hình lập phương có $6$ mặt là các hình vuông bằng nhau nên hình khai triển gồm $6$ hình vuông bằng nhau.\\
\textbf{Đáp án của câu hỏi là C}
}
\end{ex}

\begin{ex} %[5H9N49]
Hình khai triển của một hình trụ gồm những hình nào?
\choice
{$1$ hình tròn và $1$ hình chữ nhật}
{\True $2$ hình tròn bằng nhau và $1$ hình chữ nhật}
{$2$ hình vuông và $1$ hình chữ nhật}
{$2$ hình chữ nhật và $1$ hình tròn}
\loigiai{
Hình khai triển của hình trụ gồm $2$ đáy là hai hình tròn bằng nhau và mặt xung quanh là một hình chữ nhật.\\
\textbf{Đáp án của câu hỏi là B}
}
\end{ex}

\begin{ex} %[5H9N49]
Hình hộp chữ nhật có tất cả bao nhiêu đỉnh và bao nhiêu cạnh?
\choice
{\True $8$ đỉnh, $12$ cạnh}
{$6$ đỉnh, $8$ cạnh}
{$12$ đỉnh, $8$ cạnh}
{$8$ đỉnh, $6$ cạnh}
\loigiai{
Hình hộp chữ nhật có $8$ đỉnh, $12$ cạnh và $6$ mặt.\\
\textbf{Đáp án của câu hỏi là A}
}
\end{ex}

\begin{ex} %[5H9H49]
Khi mở mặt xung quanh của một hình trụ theo đường sinh, ta nhận được một hình chữ nhật có chiều dài bằng
\choice
{Chiều cao của hình trụ}
{Bán kính của hình tròn đáy}
{Diện tích của hình tròn đáy}
{\True Chu vi của hình tròn đáy}
\loigiai{
Chiều dài của hình chữ nhật bao quanh chính bằng chu vi đáy của hình trụ.\\
\textbf{Đáp án của câu hỏi là D}
}
\end{ex}

\begin{ex} %[5H9H49]
Cho một hình khai triển gồm $6$ hình vuông có cạnh $4\text{ cm}$. Khi gấp lại thành hình lập phương thì diện tích một mặt của hình lập phương đó là
\choice
{\True $16\text{ cm}^2$}
{$8\text{ cm}^2$}
{$64\text{ cm}^2$}
{$96\text{ cm}^2$}
\loigiai{
Diện tích một mặt: $4 \times 4 = 16\text{ (cm}^2\text{)}$.\\
\textbf{Đáp án của câu hỏi là A}
}
\end{ex}

\begin{ex} %[5H9H49]
Hình khai triển của hình lập phương KHÔNG thể có dạng nào sau đây?
\choice
{$4$ hình vuông thành một hàng và $2$ hình vuông ở hai bên đối diện}
{Dạng bậc thang $2 - 2 - 2$}
{\True $6$ hình vuông xếp liên tiếp thành $1$ hàng thẳng}
{Dạng $3 - 3$}
\loigiai{
Nếu $6$ hình vuông xếp thành $1$ hàng thẳng thì khi gấp lại sẽ bị đè lên nhau, không tạo thành hình hộp kín.\\
\textbf{Đáp án của câu hỏi là C}
}
\end{ex}

\begin{ex} %[5H9V49]
Một hình trụ có bán kính đáy $3\text{ cm}$, chiều cao $8\text{ cm}$. Khi khai triển mặt xung quanh của hình trụ này ta được một hình chữ nhật có diện tích là
\choice
{$75{,}36\text{ cm}^2$}
{\True $150{,}72\text{ cm}^2$}
{$120\text{ cm}^2$}
{$226{,}08\text{ cm}^2$}
\loigiai{
Chu vi đáy hình trụ: $2 \times 3{,}14 \times 3 = 18{,}84\text{ (cm)}$.\\
Diện tích mặt xung quanh (diện tích hình chữ nhật): $18{,}84 \times 8 = 150{,}72\text{ (cm}^2\text{)}$.\\
\textbf{Đáp án của câu hỏi là B}
}
\end{ex}

\begin{ex} %[5H9V49]
Người ta cắt một tấm bìa để làm một hình hộp chữ nhật không nắp có kích thước dài $10\text{ cm}$, rộng $6\text{ cm}$, cao $4\text{ cm}$. Tổng diện tích các mặt của tấm bìa hình khai triển là
\choice
{$128\text{ cm}^2$}
{$248\text{ cm}^2$}
{$240\text{ cm}^2$}
{\True $188\text{ cm}^2$}
\loigiai{
Hình hộp không nắp gồm $1$ mặt đáy và $4$ mặt xung quanh.\\
$S_{xq} = (10 + 6) \times 2 \times 4 = 128\text{ (cm}^2\text{)}$.\\
$S_{đáy} = 10 \times 6 = 60\text{ (cm}^2\text{)}$.\\
Tổng diện tích: $128 + 60 = 188\text{ (cm}^2\text{)}$.\\
\textbf{Đáp án của câu hỏi là D}
}
\end{ex}

\begin{ex} %[5H9V49]
Một hình khai triển gồm $6$ hình chữ nhật ghép lại thành hình hộp chữ nhật có các kích thước $5\text{ cm}, 4\text{ cm}, 3\text{ cm}$. Tổng diện tích toàn phần của hình khai triển đó là
\choice
{$47\text{ cm}^2$}
{\True $94\text{ cm}^2$}
{$60\text{ cm}^2$}
{$120\text{ cm}^2$}
\loigiai{
$S_{tp} = 2 \times (5 \times 4 + 5 \times 3 + 4 \times 3) = 2 \times (20 + 15 + 12) = 94\text{ (cm}^2\text{)}$.\\
\textbf{Đáp án của câu hỏi là B}
}
\end{ex}

\begin{ex} %[5H9V49]
Cần dùng một miếng bìa hình chữ nhật có kích thước tối thiểu là bao nhiêu để cuộn tròn làm thân một ống trụ có đường kính đáy $10\text{ cm}$ và chiều dài ống là $25\text{ cm}$ (không tính mép dán)?
\choice
{$10\text{ cm} \times 25\text{ cm}$}
{$15{,}7\text{ cm} \times 25\text{ cm}$}
{$62{,}8\text{ cm} \times 25\text{ cm}$}
{\True $31{,}4\text{ cm} \times 25\text{ cm}$}
\loigiai{
Chu vi đáy: $10 \times 3{,}14 = 31{,}4\text{ (cm)}$.\\
Kích thước miếng bìa là $31{,}4\text{ cm} \times 25\text{ cm}$.\\
\textbf{Đáp án của câu hỏi là D}
}
\end{ex}

\begin{ex} %[5H9V49]
Gấp một hình khai triển gồm $6$ hình vuông cạnh $a$ thành hình lập phương có thể tích $125\text{ cm}^3$. Độ dài cạnh $a$ là
\choice
{\True $5\text{ cm}$}
{$4\text{ cm}$}
{$6\text{ cm}$}
{$25\text{ cm}$}
\loigiai{
Vì $5 \times 5 \times 5 = 125$ nên $a = 5\text{ cm}$.\\
\textbf{Đáp án của câu hỏi là A}
}
\end{ex}

\begin{ex} %[5H9C49]
Trên các mặt của một con xúc xắc hình lập phương được đánh số từ $1$ đến $6$ sao cho tổng số chấm ở hai mặt đối diện luôn bằng $7$. Khi mở xúc xắc ra thành hình khai triển chữ thập ($4$ ô thẳng đứng, $2$ ô hai bên), nếu ô chính giữa là số $4$ thì ô đối diện với nó chứa số mấy?
\choice
{$1$}
{$2$}
{\True $3$}
{$5$}
\loigiai{
Mặt đối diện với số $4$ phải có số chấm là: $7 - 4 = 3$.\\
\textbf{Đáp án của câu hỏi là C}
}
\end{ex}

\textbf{Phần 2. Trắc nghiệm đúng sai}
\setcounterpageref{ex}{0}

\begin{ex} %[5H9H49]
Xét tính đúng sai về hình khai triển của các hình khối:
\choiceTF
{\True Hình khai triển của hình lập phương có đúng $6$ mặt hình vuông bằng nhau.}
{Hình khai triển của hình trụ gồm $1$ hình chữ nhật và $1$ hình tròn.}
{\True Khi cuộn mặt xung quanh của hình trụ, chiều dài hình chữ nhật bằng chu vi đáy.}
{\True Hình hộp chữ nhật có $3$ cặp mặt đối diện song song và bằng nhau.}
\loigiai{
\begin{itemize}
    \item Hình lập phương có 6 mặt vuông. Đúng.
    \item Hình trụ gồm 1 hình chữ nhật và 2 hình tròn đáy bằng nhau. Sai.
    \item Chiều dài thân trụ bằng chu vi đáy. Đúng.
    \item Hình hộp chữ nhật có 3 cặp mặt đối diện bằng nhau. Đúng.
\end{itemize}
\begin{itemize}
    \item a) \textbf{Đúng}.
    \item b) \textbf{Sai}.
    \item c) \textbf{Đúng}.
    \item d) \textbf{Đúng}.
\end{itemize}
}
\end{ex}

\begin{ex} %[5H9V49]
Một hình hộp chữ nhật có kích thước dài $8\text{ cm}$, rộng $5\text{ cm}$, cao $3\text{ cm}$. Khi trải phẳng các mặt thành hình khai triển, xét tính đúng sai:
\choiceTF
{\True Diện tích hai mặt đáy là $80\text{ cm}^2$.}
{\True Diện tích $4$ mặt bên xung quanh là $78\text{ cm}^2$.}
{Tổng diện tích toàn bộ hình khai triển là $160\text{ cm}^2$.}
{\True Tổng diện tích toàn bộ hình khai triển là $158\text{ cm}^2$.}
\loigiai{
\begin{itemize}
    \item Diện tích 2 mặt đáy: $2 \times 8 \times 5 = 80\text{ cm}^2$. Đúng.
    \item Diện tích 4 mặt xung quanh: $(8 + 5) \times 2 \times 3 = 78\text{ cm}^2$. Đúng.
    \item Tổng diện tích toàn bộ: $80 + 78 = 158\text{ cm}^2 \neq 160\text{ cm}^2$. Ý 3 Sai, Ý 4 Đúng.
\end{itemize}
\begin{itemize}
    \item a) \textbf{Đúng}.
    \item b) \textbf{Đúng}.
    \item c) \textbf{Sai}.
    \item d) \textbf{Đúng}.
\end{itemize}
}
\end{ex}

\textbf{Phần 3. Trả lời ngắn}
\setcounterpageref{ex}{0}

\begin{ex} %[5H9H49]
Hình lập phương có bao nhiêu đỉnh?
\shortans{8}
\loigiai{
\textbf{Đáp số:} $8$\
Hình lập phương có $8$ đỉnh.
}
\end{ex}

\begin{ex} %[5H9H49]
Hình khai triển của một hình lập phương gồm $6$ hình vuông, mỗi hình vuông có diện tích $25\text{ cm}^2$. Tính tổng diện tích của hình khai triển đó theo xăng-ti-mét vuông.
\shortans{150}
\loigiai{
\textbf{Đáp số:} $150$\\
$25 \times 6 = 150\text{ cm}^2$.
}
\end{ex}

\begin{ex} %[5H9V49]
Một lon sữa hình trụ có đường kính đáy $8\text{ cm}$, chiều cao $12\text{ cm}$. Nhãn giấy bọc xung quanh lon có dạng hình chữ nhật. Tính chiều dài của nhãn giấy đó theo xăng-ti-mét (lấy $\pi = 3{,}14$).
\shortans{25,12}
\loigiai{
\textbf{Đáp số:} $25{,}12$\\
Chiều dài nhãn = Chu vi đáy: $8 \times 3{,}14 = 25{,}12\text{ cm}$.
}
\end{ex}

\begin{ex} %[5H9V49]
Tính tổng diện tích các mặt của hình khai triển một hình hộp chữ nhật có chiều dài $6\text{ cm}$, chiều rộng $4\text{ cm}$, chiều cao $5\text{ cm}$ theo xăng-ti-mét vuông.
\shortans{148}
\loigiai{
\textbf{Đáp số:} $148$\\
$S_{tp} = 2 \times (6 \times 4 + 6 \times 5 + 4 \times 5) = 2 \times (24 + 30 + 20) = 148\text{ cm}^2$.
}
\end{ex}

\begin{ex} %[5H9V49]
Một hình khai triển của hình lập phương có chu vi ngoài là $56\text{ cm}$ (gồm $14$ đoạn cạnh của các hình vuông đơn vị). Tính cạnh của hình lập phương đó theo xăng-ti-mét.
\shortans{4}
\loigiai{
\textbf{Đáp số:} $4$\\
Cạnh hình lập phương: $56 : 14 = 4\text{ cm}$.
}
\end{ex}

\begin{ex} %[5H9C49]
Người ta cắt $4$ góc của một tấm bìa hình vuông cạnh $20\text{ cm}$, mỗi góc cắt đi một hình vuông nhỏ cạnh $3\text{ cm}$ rồi gập mép lên thành một chiếc hộp không nắp. Tính diện tích mặt ngoài của chiếc hộp đó theo xăng-ti-mét vuông.
\shortans{364}
\loigiai{
\textbf{Đáp số:} $364$\\
Diện tích tấm bìa ban đầu: $20 \times 20 = 400\text{ cm}^2$.\\
Diện tích $4$ góc cắt đi: $4 \times (3 \times 3) = 36\text{ cm}^2$.\\
Diện tích mặt ngoài của hộp: $400 - 36 = 364\text{ cm}^2$.
}
\end{ex}
